\subsection{Total variation between two evolving probabilities}
\label{subsec:parabolic-total-variation-growth}\index{Total variation!parabolic initial growth}

The entropies of \cref{subsec:hermitian-parabolic-entropy} compare one evolving probability with the fixed target. Two solutions of the flow with different initial potentials can nevertheless move apart. On a flat torus, even with a closed background, their total-variation distance can increase at the initial time, while no convex entropy relative to the target increases. We show this by computing the initial derivatives explicitly. By a \emph{logarithmic Monge--Amp\`ere flow with target} $\nu$ we mean a smooth solution of the equation in \eqref{eq:hermitian-parabolic-entropy-flow}, with an arbitrary smooth initial potential.

For this example, let $E=\mathbb C/(\mathbb Z+\mathrm i\mathbb Z)$ and $X=E^2$, with covering coordinates $z_j=x_j+\mathrm i y_j$. Write
\[
 \eta_j=\frac{\mathrm i}2\mathrm dz_j\wedge\mathrm d\bar z_j,
 \qquad \kappa=\eta_1+\eta_2,
 \qquad \mu=\eta_1\wedge\eta_2=\frac{\kappa^2}2.
\]
Thus $\mu$ is a probability volume form and $\int_X\kappa^2=2$. Put
\[
 c_j=\cos(2\pi x_j),\qquad S=c_1+c_2,\qquad P=c_1c_2.
\]

\begin{lemma}\label{lem:torus-total-variation-averages}
The averages of $|S|$ and $P|S|$ are
\[
 \int_X|S|\mu=\frac8{\pi^2},
 \qquad \int_XP|S|\mu=\frac8{3\pi^2}.
\]
\end{lemma}
\begin{proof}
The integrands do not depend on $y_1,y_2$. The degree-two map of real tori $(a,b)\mapsto(a+b,a-b)$ preserves normalized Haar averages. With the angles $A=2\pi a$ and $B=2\pi b$, it gives
\[
 S=2\cos A\cos B,\qquad
 P=\frac{\cos(2A)+\cos(2B)}2.
\]
The one-variable averages of $|\cos A|$ and $|\cos A|^3$ are $2/\pi$ and $4/(3\pi)$, respectively, so the average of $\cos(2A)|\cos A|$ is $2/(3\pi)$. Multiplying the averages over the two independent angles gives both formulas.
\end{proof}

For a fixed real number $K>3$, define the smooth target probability $\nu_K$ by
\[
 Z_K=\int_X\mathrm e^{KP}\mu,
 \qquad f_K=Z_K^{-1}\mathrm e^{KP},
 \qquad \nu_K=f_K\mu.
\]
For $0<\varepsilon<1/\pi$, the initial potentials $\varepsilon S$ and zero give positive forms. Set
\[
 A_\varepsilon=1-\pi\varepsilon c_1,
 \qquad B_\varepsilon=1-\pi\varepsilon c_2,
 \qquad r_\varepsilon=A_\varepsilon B_\varepsilon,
\]
so that $\kappa+\ddc(\varepsilon S)=A_\varepsilon\eta_1+B_\varepsilon\eta_2$ and its normalized Monge--Amp\`ere probability is $r_\varepsilon\mu$.

Consider the short smooth paths
\begin{equation}\label{eq:torus-parabolic-initial-paths}
 U_\varepsilon(t)=\varepsilon S+t\log\frac{r_\varepsilon}{f_K},
 \qquad V(t)=-t\log f_K.
\end{equation}
They remain positive for sufficiently small $|t|$, and at $t=0$ their velocities are exactly those prescribed by the logarithmic Monge--Amp\`ere flow with target $\nu_K$. Their normalized probabilities are
\[
 \mu_\varepsilon(t)=\frac{(\kappa+\ddc U_\varepsilon(t))^2}2,
 \qquad \mu_0(t)=\frac{(\kappa+\ddc V(t))^2}2,
\]
which have mass one because the background $\kappa$ is closed. Define their total-variation distance, with the convention for probabilities, by
\[
 D_\varepsilon(t)
 =\frac12\int_X|\mu_\varepsilon(t)-\mu_0(t)|.
\]

\Needspace{12\baselineskip}
\begin{proposition}\label{prop:parabolic-total-variation-growth}\index{Complex torus!parabolic total-variation example}
For the paths \eqref{eq:torus-parabolic-initial-paths},
\[
 D'_\varepsilon(0)
 =4\varepsilon\left(\frac K3-1\right)+o(\varepsilon)
 \quad\text{as }\varepsilon\downarrow0.
\]
In particular, this derivative is positive for sufficiently small $\varepsilon>0$. Any two logarithmic Monge--Amp\`ere flows with target $\nu_K$ and these initial potentials have the same positive initial derivative of their total-variation distance.
\end{proposition}
\begin{proof}
For a real function $h$ depending only on $x_1,x_2$, the coefficient matrix of $\ddc h$ with respect to the forms $(\mathrm i/2)\mathrm dz_j\wedge\mathrm d\bar z_k$ is $(4\pi)^{-1}(h_{x_jx_k})$. Put $R_\varepsilon(t)=\mathrm d\mu_\varepsilon(t)/\mathrm d\mu$ and $R_0(t)=\mathrm d\mu_0(t)/\mathrm d\mu$. Differentiating the determinant at the diagonal initial matrices gives
\[
\begin{split}
 \left.\partial_t(R_\varepsilon-R_0)\right|_{t=0}
 =\frac1{4\pi}\bigl(&B_\varepsilon\partial_{x_1}^2\log A_\varepsilon
             +A_\varepsilon\partial_{x_2}^2\log B_\varepsilon\\
 &-K(B_\varepsilon-1)\partial_{x_1}^2P
             -K(A_\varepsilon-1)\partial_{x_2}^2P\bigr).
\end{split}
\]
Since $\partial_{x_j}^2P=-4\pi^2P$, Taylor expansion in $\varepsilon$ yields the uniform estimate
\[
 \left.\partial_t(R_\varepsilon-R_0)\right|_{t=0}
 =\pi^2\varepsilon S(1-KP)+O(\varepsilon^2).
\]
At $t=0$, the difference of the densities is
\[
 q_\varepsilon=r_\varepsilon-1
 =-\pi\varepsilon S+\pi^2\varepsilon^2P.
\]
For each fixed $x_2$, its coefficient of $c_1$ is $-\pi\varepsilon(1-\pi\varepsilon c_2)\ne0$, so as a function of $x_1$ it has at most two zeros in a period, and its zero set has measure zero. Differentiation under the integral sign, justified by the uniform bound on the derivatives of the smooth densities, gives
\[
 D'_\varepsilon(0)=\frac12\int_X\operatorname{sgn}(q_\varepsilon)
       \left.\partial_t(R_\varepsilon-R_0)\right|_{t=0}\mu.
\]
The set $\{S=0\}$ also has measure zero, and $\operatorname{sgn}(q_\varepsilon)\to-\operatorname{sgn}(S)$ off this set. Dominated convergence and \cref{lem:torus-total-variation-averages} give
\[
 \lim_{\varepsilon\downarrow0}\frac{D'_\varepsilon(0)}\varepsilon
 =\frac{\pi^2}2\int_X|S|(KP-1)\mu
 =4\left(\frac K3-1\right).
\]
Every flow with the stated initial data has the velocities \eqref{eq:torus-parabolic-initial-paths} at $t=0$, so its initial density derivatives, and hence the derivative of the distance, coincide with this computation.
\end{proof}

\begin{corollary}\label{cor:individual-entropy-pair-growth}\index{Monge--Ampere equation@Monge--Amp\`ere equation!pairwise parabolic variation}
For these paths, the initial derivative of every $C^2$ convex entropy of each probability relative to $\nu_K$ is nonpositive, while the total-variation distance between the two probabilities has positive initial derivative for sufficiently small $\varepsilon>0$.
\end{corollary}
\begin{proof}
The initial velocities are those of the logarithmic flow, and both initial forms are closed. The chain-rule calculation in the proof of \cref{prop:hermitian-parabolic-entropy-identity} therefore applies at $t=0$: its background term vanishes and its gradient term is nonpositive. The preceding proposition gives the assertion about the pair.
\end{proof}

Since only initial velocities enter, the example shows, without any existence or convergence theorem for the flow, that the logarithmic Monge--Amp\`ere flow need not contract the total-variation distance infinitesimally, even though no convex entropy relative to the target increases at the initial time.
